\subsection{Feature Ablation Analysis}
\label{sec:feature_ablation}

Kuantifikasi perubahan akurasi klasifikasi dari skema domain tunggal menuju konfigurasi ganda dan tiga domain dirangkum pada Table~\ref{tab:VII}, Table~\ref{tab:VIII}, Table~\ref{tab:IX}, dan Table~\ref{tab:X}. Pada model SVM, transisi representasi dari Face (90.83\%) ke Emotion $\oplus$ Face (93.29\%) meningkatkan akurasi sebesar 2.46\% (0.0246). Penambahan domain ketiga pada konfigurasi tri-domain Face $\oplus$ Emotion $\oplus$ Age menghasilkan peningkatan performa lebih lanjut menjadi 93.70\%, yang mencerminkan peningkatan kumulatif sebesar 2.87\% (0.0287) di atas fitur tunggal Face. Pola peningkatan progresif ini mengindikasikan kehadiran informasi diskriminatif tambahan dari setiap domain representasi visual yang digabungkan, meskipun mekanisme interaksi antardimensi fitur laten tidak dapat dipastikan hanya dari hasil pengujian empiris.

Analisis kontribusi representasi domain usia (Age) membedah peranan informasi laten penuaan dalam sistem klasifikasi demografis. Pada seluruh pengklasifikasi yang dievaluasi, Age merupakan konfigurasi single-domain dengan capaian akurasi terendah, seperti tercatat sebesar 87.64\% pada SVM dan 69.63\% pada GNB. Namun demikian, penggabungan fitur Age dengan fitur biometrik wajah pada skema dual-domain Face $\oplus$ Age menghasilkan akurasi 92.55\%, yang memberikan peningkatan sebesar 1.72\% (0.0172) di atas performa fitur Face murni (90.83\%) pada model SVM. Hasil empiris tersebut menunjukkan kemungkinan adanya informasi diskriminatif tambahan dari age-associated representations yang melengkapi profil morfologi wajah, tanpa mengklaim bahwa kontribusi informasi tersebut bersifat ortogonal atau komplementer mutlak secara teoretis.
